% year: תשע"ט
% semester: ב
% moed: א
% number: 3ב
% points: 10
% categories: taylor
% source: exams/מבחנים/תשעט/חודיסמן מועד א תשעט.pdf

\begin{question}
(10 נק') רשמו פולינום מקלורן מסדר 2 עבור הפונקציה $f(x)=\cos(2x)-\arctan(x+1)$.
\end{question}

\begin{hint}
פתחו כל מחובר בנפרד: ל-$\cos(2x)$ השתמשו בפיתוח הידוע של $\cos t$, ול-$\arctan(x+1)$ חשבו את $g(0),g'(0),g''(0)$.
\end{hint}

\begin{solution}
פולינום מקלורן מסדר 2: $P_2(x)=f(0)+f'(0)x+\frac{f''(0)}{2}x^2$. נחשב בנפרד לכל מחובר.

עבור $\cos(2x)$: לפי הפיתוח $\cos t=1-\frac{t^2}{2}+o(t^2)$ נקבל $\cos 2x=1-2x^2+o(x^2)$.

עבור $g(x)=\arctan(x+1)$:
\[ g(0)=\arctan1=\frac\pi4,\quad g'(x)=\frac{1}{1+(x+1)^2}\Rightarrow g'(0)=\frac12,\quad g''(x)=-\frac{2(x+1)}{\big(1+(x+1)^2\big)^2}\Rightarrow g''(0)=-\frac24=-\frac12. \]
לכן $\arctan(x+1)=\frac\pi4+\frac x2-\frac{x^2}{4}+o(x^2)$. בחיסור:
\[ P_2(x)=\left(1-\frac\pi4\right)-\frac12x-\frac74x^2. \]
\end{solution}
